\scanpage{24}
qui contient $\hat{\mathcal E}$ comme sous-extension. On désigne par $\Gamma$ le commutant de $\unsure{\mathfrak S_3|\mathcal S_3}$ dans $\widehat{\widehat\Gamma}$ \inserted{(\unsure{qui centralise} $\unsure{\tau}$)}, et par $E$ l'extension correspondante de $\Gamma$ par $\hat\pi$. Notons déjà que
\[
\left\{\begin{aligned}
\Gamma\cap\deleted{\unsure{\mathfrak S_3}}\Gamma_{\unsure{0,3}}&=\{1,\tau\}\\
\Gamma\cap\unsure{\mathfrak S_3|\mathcal S_3}&=\{1\}.
\end{aligned}\right.
\]
\margin{\deleted{\uline{NB}}}
\margin{\uline{NB}
\[
\diagram{\begin{tikzcd}
\unsure{\mathfrak S_3|\mathcal S_3}\times\widehat{\widehat\Gamma}^!\arrow[d,equal]&&\\
\widehat{\widehat\Gamma}\arrow[r,phantom,"\supset"]&\unsure{\mathfrak S_3|\mathcal S_3}\times\Gamma&
\end{tikzcd}}
\]
--- j'ignore si cette inclusion est une égalité, i.e. si $\Gamma\subset\widehat{\widehat\Gamma}^!$ est une égalité.}
On veut donner d'abord une description plus symétrique des élts de $\widehat{\widehat{\mathcal E}}^!$, sans faire jouer un rôle particulier à $L_0$ engendré par $l_0$ $(L_0=\hat\kappa_0(\hat{\mathbb Z}))$. Un $u\in\unsure{\hat{\mathcal E}|\widehat{\widehat{\mathcal E}}}$ est \inserted{pour être} défini par un multiplicateur
\[
(19)\qquad p=\chi(u)\in\hat{\mathbb Z}^{*}
\]
et des
\[
(20)\qquad g_i\in\hat\pi\qquad i=0,1,\infty
\]
pour
\[
(21)\qquad u(l_i)=\operatorname{int}(g_i^{\unsure{\deleted{-1}|u}})l_i^p
\]
les données (19), (20) étant liées par la relation
\[
(22)\qquad(\operatorname{int}(g_\infty^{\unsure{\deleted{-1}|u}})l_\infty^p)
(\operatorname{int}g_1^{\unsure{\deleted{-1}|u}}\cdot l_1^p)
(\operatorname{int}(g_0^{\unsure{\deleted{-1}|u}})l_0^p)=1
\]
et par le fait que les $u(l_i)$ définis par (21) \uline{engendrent} $\hat\pi$ topologiquement. Pour $u$ fixé, $p$
\mspage{23}
